\chapter{Grundlagen}
\label{ch:grundlagen}

Dieses Kapitel führt die Begriffe und Verfahren ein, die für die Arbeit
vorausgesetzt werden. Es reicht vom Reinforcement Learning über die
Grundbegriffe der Netzwerksicherheit bis zu den Arbeiten, an die die
Untersuchung anschließt. Wie diese Grundlagen für die
eigenen Versuche verwendet werden, behandelt \Cref{ch:versuchsaufbau}.

Als Beispielszenario dient die Abwehr von Angriffen in einem
Unternehmensnetz. Es begleitet die Darstellung von den allgemeinen
Lernverfahren bis zu den sicherheitstechnischen Begriffen und verbindet so
beide Hälften des Kapitels.

\section{Reinforcement Learning}
\label{sec:rl}

Dieser Abschnitt entwickelt das \ac{RL} schrittweise von seiner Einordnung im
maschinellen Lernen bis zu den Grundlagen, die der weitere Text voraussetzt.

\subsection{Einordnung in das maschinelle Lernen}

Maschinelles Lernen umfasst Verfahren, die eine Aufgabe nicht über eine
ausformulierte Regel lösen, sondern ihr Verhalten aus Erfahrung ableiten. Welche
Art von Rückmeldung diese Erfahrung liefert, trennt die Teilgebiete. Neben dem
überwachten und dem unüberwachten Lernen führen Sutton und Barto entlang dieses
Kriteriums das \ac{RL} als dritte Form
an~\cite[Abschn.~1.1]{suttonbarto2018}. Für diese dritte Form ist im deutschen Sprachgebrauch keine Bezeichnung
gebräuchlich. Sie wird deshalb durchgehend
englisch benannt.

Beim \emph{überwachten Lernen} ist die Rückmeldung \emph{instruktiv}. Ein
Klassifikator erhält Mitschnitte des Netzwerkverkehrs, die vorab als gutartig
oder bösartig eingestuft wurden. Weil zu jeder Verbindung die korrekte
Einstufung bekannt ist, lässt sich der Fehler des Modells unmittelbar berechnen.
Beim \emph{unüberwachten Lernen} entfällt jede Bewertung. Derselbe Klassifikator
erhält dieselben Mitschnitte ohne Einstufung und fasst Verbindungen mit
ähnlichem Muster zu Gruppen zusammen, ohne zu wissen, ob eine dieser Gruppen
einen Angriff darstellt.

Beim \ac{RL} ist die Rückmeldung \emph{evaluativ}. Hier lässt sich das Beispiel
des Klassifikators nicht mehr fortführen, denn dieser ordnet vorliegende Daten
ein, ohne in das Geschehen einzugreifen. An seine Stelle tritt ein Agent,
der handelt und für jede Handlung eine Bewertung erhält, aber keinen Hinweis
darauf, welche Handlung besser gewesen wäre. Die Rückmeldung bewertet also, ohne
anzuleiten~\cite[S.~25]{suttonbarto2018}. Im vorliegenden Beispiel entspricht das einem Verteidiger, der auf
die Lage im Netz mit einer Gegenmaßnahme reagiert und allein daran gemessen
wird, wie sich diese auswirkt.

Ein zweiter Unterschied betrifft nicht die Rückmeldung, sondern die Herkunft
der Daten. In den ersten beiden Fällen liegen die Mitschnitte vor und sind vom
Modell unabhängig; man könnte dasselbe Modell zweimal auf denselben Daten
trainieren. Im \ac{RL} entsteht die Erfahrung dagegen erst durch das Handeln
des Agenten: Seine Maßnahmen bestimmen, welche Situationen er anschließend
überhaupt zu sehen bekommt. Sichert ein Verteidiger einen Netzbereich nicht
ab, so kann er auch nicht die Zustände beobachten, welche aus dieser Maßnahme
hervorgegangen wären.

In beiden zuletzt genannten Fällen sind die Mitschnitte nicht eingestuft, was
\ac{RL} als Sonderfall des unüberwachten Lernens erscheinen lässt. Sutton und
Barto widersprechen dem ausdrücklich~\cite[Abschn.~1.1]{suttonbarto2018}:
Maßgeblich ist nicht die fehlende Einstufung, sondern das Ziel. Unüberwachtes
Lernen sucht verborgene Struktur, \ac{RL} maximiert ein Reward-Signal.

\subsection{Das Grundprinzip des Reinforcement Learnings}
\label{sec:rl-grundprinzip}

Ein \ac{RL}-Aufbau besteht aus zwei Teilen, dem \emph{Agenten} und der
\emph{Umgebung}, die in einer festen Schleife miteinander interagieren. Zu jedem
diskreten Zeitschritt $t$ beobachtet der Agent den \emph{Zustand} $s_t$ der
Umgebung, wählt daraufhin eine \emph{Aktion} $a_t$ und erhält von der Umgebung
zwei Daten zurück: den neuen Zustand $s_{t+1}$ und ein Reward-Signal $R_{t+1}$,
also eine einzelne Zahl, die diesen Schritt bewertet.

Wie lange diese Schleife läuft, hängt von der Aufgabe ab. Sutton und Barto
unterscheiden \emph{fortlaufende} Aufgaben, die kein natürliches Ende kennen,
von \emph{episodischen}, die einen Endzustand
besitzen~\cite[Abschn.~3.3]{suttonbarto2018}; welche Folgen das hat, zeigt
\Cref{sec:mdp}. Die vorliegende Arbeit betrachtet ausschließlich den zweiten
Fall. Ein Durchlauf vom Anfangszustand bis zum Endzustand heißt \emph{Episode};
ein Training umfasst eine Vielzahl solcher Episoden. Die Umgebung wird dabei
zu Beginn jeder Episode zurückgesetzt, die gelernten Parameter des Agenten
bleiben jedoch erhalten.

Was der Agent dabei nicht erfährt, ist ebenso wesentlich. Ob eine andere Wahl zu einem höheren Reward
geführt hätte, muss er durch wiederholtes Ausprobieren erschließen.
Sein Ziel ist außerdem nicht der einzelne Reward, sondern deren Summe über die
gesamte Episode. Das schließt den Fall ein, dass eine Aktion mit geringem unmittelbarem
Reward erst die Voraussetzung für spätere, höhere Rewards schafft.

Dahinter steht eine Annahme, die Sutton und Barto als \emph{Reward-Hypothese}
benennen. Jedes Ziel lasse sich als Maximierung des Erwartungswerts einer
einzelnen skalaren Größe auffassen, der Summe der erhaltenen
Rewards~\cite[Abschn.~3.2]{suttonbarto2018}. Ob eine konkrete
Reward-Funktion das beabsichtigte Ziel tatsächlich abbildet, ist damit
nicht sichergestellt; \Cref{sec:rl-schwierigkeiten} greift dies wieder auf.

Vier Bestandteile setzen das Verfahren zusammen. Die \emph{Policy} legt das Verhalten des
Agenten fest, indem sie jedem Zustand die dort zu wählende Aktion zuordnet. Das
\emph{Reward-Signal} gibt das Ziel vor und bewertet jeden Schritt einzeln. Die
\emph{Wertfunktion} schätzt demgegenüber die Summe der Rewards, die von
einem Zustand aus noch zu erwarten sind, und blickt damit über den
einzelnen Schritt hinaus.
Hinzu kommt gegebenenfalls ein \emph{Modell} der Umgebung, das vorhersagt, wie
diese auf eine Aktion reagieren
wird~\cite[Abschn.~1.3]{suttonbarto2018}. Die folgenden Abschnitte führen diese
Bestandteile der Reihe nach ein.

\subsection{Verzögerte Belohnung}
\label{sec:credit}

Eine zentrale Herausforderung des Verfahrens besteht darin, dass ein Reward und
die für ihn ursächliche Aktion zeitlich weit auseinanderliegen können. Eine
Aktion kann unmittelbar keinen Reward erbringen und dennoch die Voraussetzung
für einen Erfolg schaffen, der erst viele Schritte später eintritt. Der Reward
fällt am Ende dieser Kette an, ohne dass erkennbar wäre, welchem der
vorangegangenen Schritte er zuzurechnen ist.

Für das vorliegende Beispiel: Übernimmt ein Angreifer einen unscheinbaren
Rechner, bringt ihm das kaum Punkte. Findet er dort jedoch Zugangsdaten, die ihm
zwanzig Schritte später den Zugriff auf ein wertvolles System eröffnen, war
genau dieser frühe Schritt ausschlaggebend.

Minsky benannte diese Schwierigkeit bereits 1961 als
\emph{Credit-Assignment-Problem}~\cite{minsky1961}. Es ist der Hauptgrund
dafür, dass \ac{RL} sehr viele Durchläufe benötigt: Die Bewertung eines frühen
Schrittes lässt sich erst festigen, nachdem sein Nutzen über viele Episoden
hinweg wiederholt sichtbar geworden ist. Sämtliche im Folgenden eingeführten
Größen sind Werkzeuge, die dieses Problem adressieren.

\subsection{Der Markov-Entscheidungsprozess}
\label{sec:mdp}

Um das Zusammenspiel aus \Cref{sec:rl-grundprinzip} präzise zu beschreiben, wird
es als \ac{MDP} formalisiert. Hinter dem Begriff verbirgt sich eine Sammlung von
fünf Angaben, die jede \ac{RL}-Aufgabe ohnehin festlegen muss:

\begin{enumerate}
  \item die Menge aller möglichen Zustände $\mathcal{S}$,
  \item die Menge aller Aktionen $\mathcal{A}$,
  \item die Übergangsfunktion $P(s' \mid s, a)$, also die Wahrscheinlichkeit,
        nach Aktion $a$ im Zustand $s$ im Zustand $s'$ zu landen,
  \item die Reward-Funktion $R(s, a)$, welche die unmittelbare Bewertung
        liefert,
  \item den Diskontfaktor $\gamma$ mit $0 \le \gamma \le 1$, der festlegt, wie
        stark künftige Rewards gegenüber sofortigen zählen.
\end{enumerate}

Benannt ist das Modell nach dem Mathematiker Andrei Markov, auf dessen Arbeiten
zu stochastischen Prozessen die zugrunde liegende Annahme zurückgeht: Der
Zustand enthält bereits alle Information, die für die Zukunft von Bedeutung ist.
Wie der Agent dorthin gelangt ist, spielt keine Rolle. Beim Schach trifft das
zu, denn für den nächsten Zug zählt allein die aktuelle Stellung, nicht die
Zugfolge, die zu ihr geführt hat. Die Annahme vereinfacht das Problem erheblich,
weil der Agent keine Vorgeschichte mitführen
muss~\cite[Abschn.~3.1]{suttonbarto2018}.

Beide Randwerte des Diskontfaktors sind zulässig, führen aber zu Sonderfällen.
Bei $\gamma = 0$ zählt ausschließlich der unmittelbare Reward; der Agent handelt
vollständig kurzsichtig und kann verzögerte Belohnungen aus \Cref{sec:credit}
grundsätzlich nicht mehr erkennen. Bei $\gamma = 1$ zählt jeder künftige Reward
mit vollem Gewicht. Zulässig ist das nur bei episodischen Aufgaben im Sinne von
\Cref{sec:rl-grundprinzip}, bei denen jeder Durchlauf sicher endet. Eine
fortlaufende Aufgabe liefert dagegen unendlich viele Rewards, deren Summe ohne
Abwertung unbeschränkt anwachsen kann; das Ziel des Agenten wäre dann nicht mehr
definiert. Ein Wert knapp darunter genügt bereits: Die Beiträge nehmen dann
geometrisch ab, und die Summe konvergiert. 

Die fünf Angaben beschreiben die Aufgabe. Ob der Agent die Übergangsfunktion $P$
und die Reward-Funktion $R$ seinerseits kennt, ist damit noch nicht gesagt, und
davon hängt ab, wie er lernen kann~\cite[Abschn.~1.3]{suttonbarto2018}. Beim Schach
kennt er beide, denn die Regeln bestimmen, welche Stellung auf einen Zug folgt
und ob sie gewonnen ist; er kann Züge daher durchspielen, ohne sie auszuführen.
Solche Verfahren heißen \emph{modellbasiert}. Ein Verteidiger im
Unternehmensnetz hat diese Möglichkeit nicht: Sperrt er eine Verbindung, weiß er
weder, ob der Angreifer sie überhaupt nutzen wollte, noch ob dieser auf einen
anderen Weg ausweicht. Er erfährt erst hinterher, was eingetreten ist, und
Verfahren, die allein damit arbeiten, heißen \emph{modellfrei}.

Die Agenten dieser Arbeit lernen modellfrei, und eine Alternative besteht nicht:
Eine Übergangsfunktion für ein Netz unter Angriff aufzustellen setzte voraus,
den Ausgang jedes möglichen Angriffsschritts vorab zu kennen. Was sich nicht
vorausrechnen lässt, muss ausprobiert werden; daraus folgt der hohe Bedarf an
Erfahrung, den \Cref{sec:rl-schwierigkeiten} aufgreift.

\subsection{Policy, Return und Diskontierung}
\label{sec:policy}

Das Verhalten des Agenten beschreibt die \emph{Policy} $\pi(a \mid s)$. Sie ist
eine bedingte Wahrscheinlichkeit: $\pi(a \mid s)$ gibt an, mit welcher
Wahrscheinlichkeit der Agent die Aktion $a$ ausführt, wenn er sich im Zustand
$s$ befindet. Für jeden festen Zustand ergibt sie damit eine
Wahrscheinlichkeitsverteilung über alle Aktionen, deren Werte sich zu eins
summieren. Die Policy ist genau das, was im Training verändert wird.

Ihr Ziel ist nicht der größtmögliche unmittelbare Reward. Maßgeblich ist der
\emph{Return} $G_t$ zum Zeitpunkt $t$, also die Summe aller noch folgenden
Rewards $R_{t+k+1}$, wobei spätere Beiträge mit dem Diskontfaktor $\gamma$
abgewertet werden:
%
\begin{equation}
  G_t = \sum_{k=0}^{\infty} \gamma^{k} R_{t+k+1}.
  \label{eq:return}
\end{equation}
%
Der Laufindex $k$ zählt dabei die künftigen Schritte, sodass der Faktor
$\gamma^{k}$ einen Reward umso stärker abwertet, je weiter er in der Zukunft
liegt. Der um eins verschobene Index folgt der Zählweise aus
\Cref{sec:rl-grundprinzip}: Ein Reward wird nach dem Zeitpunkt benannt, zu dem
er \emph{eintrifft}, und er trifft stets einen Schritt nach der Aktion ein, die
ihn ausgelöst hat. Für $k = 0$ steht dort deshalb $R_{t+1}$, also der Reward
für die zum Zeitpunkt $t$ gewählte Aktion.

\subsection{Wert- und Vorteilsfunktion}
\label{sec:wertfunktion}

Um Zustände und Aktionen vergleichbar zu machen, dienen zwei Größen. Die
\emph{Aktionswertfunktion} $Q^{\pi}(s,a)$ gibt den erwarteten Return an, wenn
der Agent im Zustand $s$ zunächst die Aktion $a$ ausführt und erst danach wieder
der Policy $\pi$ folgt. Die \emph{Zustandswertfunktion} $V^{\pi}(s)$ beantwortet
dieselbe Frage ohne Festlegung auf eine bestimmte Aktion: Sie ist der
Erwartungswert von $Q^{\pi}(s,a)$ über die Aktionen, die die Policy im Zustand
$s$ wählt. Bei endlich vielen Aktionen ist das die mit den
Auswahlwahrscheinlichkeiten gewichtete Summe:
%
\begin{equation}
  V^{\pi}(s)
    = \mathbb{E}_{a \sim \pi(\cdot \mid s)}\!\left[Q^{\pi}(s,a)\right]
    = \sum_{a \in \mathcal{A}} \pi(a \mid s)\, Q^{\pi}(s,a).
  \label{eq:v-aus-q}
\end{equation}
%
Dabei bezeichnet $\mathbb{E}$ den Erwartungswert, und die Schreibweise
$a \sim \pi(\cdot \mid s)$ besagt, dass die Aktion $a$ gemäß der Policy im
Zustand $s$ gezogen wird. Auch $V^{\pi}$ setzt Handeln voraus, der Unterschied
zu $Q^{\pi}$ liegt allein darin, dass über die Aktionen gemittelt wird, statt
eine einzelne festzulegen. Der Zustandswert steht damit nicht für eine bestimmte
Aktion, etwa die wahrscheinlichste, sondern für den Erwartungswert dessen, was
die Policy in $s$ tun würde.

Aus der Differenz beider ergibt sich die \emph{Vorteilsfunktion}, benannt nach
dem englischen \emph{advantage} und nicht zu verwechseln mit der Aktionsmenge
$\mathcal{A}$:
%
\begin{equation}
  A^{\pi}(s,a) = Q^{\pi}(s,a) - V^{\pi}(s).
  \label{eq:advantage}
\end{equation}
%
Sie beantwortet die für das Training entscheidende Frage: War diese Aktion
besser als das, was die Policy in dieser Lage üblicherweise tut? Ein positiver
Wert bedeutet besser als der eigene Erwartungswert, ein negativer schlechter.
Lernverfahren verschieben die Policy in Richtung der Aktionen mit positivem
Vorteil; wie das im Einzelnen geschieht, behandelt \Cref{sec:policyupdate}.

Ein Zahlenbeispiel verdeutlicht das Zusammenspiel. Stünden dem Agenten in einer
Lage drei Aktionen mit den Aktionswerten 80, 10 und 0 offen und wählte die
Policy sie mit 70, 20 und 10 Prozent, ergäbe sich nach \Cref{eq:v-aus-q} ein
Zustandswert von $0{,}7 \cdot 80 + 0{,}2 \cdot 10 + 0{,}1 \cdot 0 = 58$. Die
erste Aktion hätte damit den Vorteil $+22$, die zweite $-48$. Steigt der Anteil
der ersten im Lauf des Trainings auf 95 Prozent, wächst der Zustandswert auf
76{,}4, und ihr Vorteil schrumpft auf $+3{,}6$.

Die Bezugsgröße wächst also mit der Policy mit. Sobald der Agent überwiegend das
Beste tut, kann nichts mehr deutlich besser als sein eigener Erwartungswert
sein, und die Anpassungen kommen zum Erliegen. Formal folgt das aus
\Cref{eq:v-aus-q}: Unter der Policy gemittelt ergibt der Vorteil stets null.

Der Vorteil misst nur den Abstand zum eigenen Erwartungswert. Ob eine Lage
insgesamt günstig oder ungünstig ist, fällt damit heraus; verglichen wird
ausschließlich innerhalb einer Lage. Unempfindlich gegen die Höhe der Rewards
ist das Verfahren deshalb aber nicht: Werden alle Rewards mit zehn
multipliziert, wächst der Vorteil um denselben Faktor und mit ihm die Größe der
Anpassung. Heraus fällt das Niveau einer einzelnen Lage, nicht der Maßstab der
Reward-Funktion.

$V^{\pi}$ und $Q^{\pi}$ sind dem Agenten allerdings nicht bekannt. Wären sie es,
ließe sich stets die Aktion mit dem höchsten Aktionswert wählen, und ein
Lernproblem bestünde nicht. Beide müssen deshalb geschätzt werden, und wie das
geschieht, behandelt \Cref{sec:drl}.

\subsection{Actor und Critic}
\label{sec:actorcritic}

Aus den bisher eingeführten Größen ergibt sich eine Aufgabenteilung. Ein Agent
muss zweierlei leisten: Er muss eine Aktion wählen, und er muss einschätzen
können, wie gut seine Lage ist. Das Erste leistet die Policy aus
\Cref{sec:policy}, das Zweite die Zustandswertfunktion aus
\Cref{sec:wertfunktion}. Beide sind bislang nur beschrieben, nicht verortet.

Eine verbreitete Familie von Verfahren weist die beiden Aufgaben deshalb
getrennten Bestandteilen zu, und das erklärt zugleich deren
Namen~\cite[Abschn.~13.5]{suttonbarto2018}. Der
\emph{Actor} hält die Policy und wählt die Aktionen; er ist der handelnde Teil.
Der \emph{Critic} schätzt den Zustandswert und bewertet im Nachhinein, ob das
Ergebnis besser oder schlechter ausfiel als erwartet. Diese Differenz ist der
Vorteil aus \Cref{eq:advantage}. Er steuert die Anpassung des Actors, während
der Critic zugleich an den beobachteten Rewards nachjustiert wird. Beide
verbessern sich also im selben Durchlauf; zu Beginn ist der Critic allerdings
ungenau, weshalb frühe Vorteilsschätzungen mit Vorsicht zu betrachten sind.

Diese Zweiteilung ist nicht zwingend. \emph{Wertbasierte} Verfahren kommen ohne
eigenen Actor aus und leiten ihr Verhalten unmittelbar aus geschätzten
Aktionswerten ab; rein \emph{policybasierte} Verfahren verzichten umgekehrt auf
den Critic und arbeiten allein mit den beobachteten
Returns~\cite[Abschn.~13.3--13.5]{suttonbarto2018}.
\emph{Actor-Critic-Verfahren} verbinden beides. Der in dieser Arbeit
eingesetzte Algorithmus zählt zu dieser Klasse.

\subsection{Exploration und Exploitation}

Ein Agent, der stets die bislang beste bekannte Aktion wählt, entdeckt nichts
Besseres mehr. Einer, der nur ausprobiert, nutzt sein Wissen nicht. Zwischen
diesen beiden Polen muss abgewogen werden, und die Abwägung ist nicht einmalig,
sondern begleitet das gesamte Training. Sutton und Barto heben diesen Konflikt
als eine Schwierigkeit hervor, die im überwachten und unüberwachten Lernen gar
nicht erst auftritt~\cite[Abschn.~1.1]{suttonbarto2018}.

Verfahren mit \emph{stochastischer} Policy lösen das implizit. Weil die Policy
nach \Cref{sec:policy} eine Wahrscheinlichkeitsverteilung ausgibt und keine
einzelne Aktion, bleibt ein Rest an Exploration erhalten, solange diese
Verteilung nicht auf eine einzige Aktion zusammenfällt. Der Agent zieht seine
Aktion in jedem Schritt neu aus der Verteilung und weicht damit gelegentlich von
seiner derzeit besten Einschätzung ab.

\section{Deep Reinforcement Learning}
\label{sec:drl}

Der vorangegangene Abschnitt beschreibt das \ac{RL} unabhängig davon, wie
Policy und Wertfunktion dargestellt werden. Dieser Abschnitt führt tiefe
neuronale Netze als solche Darstellung ein, beschreibt die daraus
entstehende Actor-Critic-Architektur und das in dieser Arbeit eingesetzte
Verfahren \ac{PPO} und schließt mit den Schwierigkeiten, die dieses
Vorgehen mit sich bringt.

\subsection{Von der Tabelle zur Funktionsapproximation}
\label{sec:funktionsapproximation}

Wie Policy und Wertfunktionen dargestellt werden, blieb bisher offen. Für kleine
Aufgaben genügt eine Tabelle mit einem Eintrag je Zustand. Die Einträge stehen
dabei unabhängig nebeneinander: Was der Agent in einem Zustand gelernt hat,
steht in genau einer Zeile und sagt über keine andere etwas aus. Bei großen
Zustandsräumen wird das zum Problem, denn dort erreicht der Agent jeden
einzelnen Zustand nur selten oder überhaupt nicht. Er müsste jeden davon eigens
erleben, um ihn beurteilen zu können.

Die \emph{Funktionsapproximation} ersetzt die Tabelle durch eine Funktion, die
den Wert aus den Merkmalen eines Zustands \emph{berechnet}, statt ihn
nachzuschlagen. Sie besitzt einstellbare Parameter, die im Training angepasst
werden; bei neuronalen Netzen heißen diese \emph{Gewichte} und werden
zusammenfassend mit $\theta$ bezeichnet. Weil ähnliche Zustände ähnliche
Merkmale aufweisen, liefert die Funktion für sie auch ähnliche Ergebnisse.
Gelerntes überträgt sich damit auf Zustände, die der Agent nie besucht hat.

Werden für diese Approximation tiefe neuronale Netze eingesetzt, spricht man von
\emph{Deep Reinforcement Learning}. Der Ansatz verbindet die Entscheidungslogik
des \ac{RL} mit der Fähigkeit neuronaler Netze, aus hochdimensionalen
Beobachtungen brauchbare Merkmale zu gewinnen~\cite[Teil~II]{suttonbarto2018}.
Auf diese Weise lernten Agenten bereits allein aus Bildpunkten und Punktestand,
Videospiele auf menschlichem Niveau zu spielen~\cite[S.~529]{mnih2015dqn}.

\subsection{Die Actor-Critic-Architektur}
\label{sec:ac-architektur}

Actor und Critic aus \Cref{sec:actorcritic} sind bislang Rollen, keine
konkreten Gebilde: Es steht fest, was sie leisten sollen, nicht aber, wodurch.
Die Funktionsapproximation beantwortet das. Beide Rollen werden durch je ein
neuronales Netz besetzt, und diese Aufteilung heißt
\emph{Actor-Critic-Architektur}~\cite[Abschn.~13.5]{suttonbarto2018}.

Der \emph{Actor} bildet den beobachteten Zustand auf eine
Wahrscheinlichkeitsverteilung über die Aktionen ab und stellt damit die Policy
dar; seine Gewichte werden mit $\theta$ bezeichnet, weshalb die Policy im
Folgenden als $\pi_{\theta}$ geschrieben wird. Der \emph{Critic} bildet
denselben Zustand auf eine einzelne Zahl ab, nämlich auf seine Schätzung des
Zustandswerts $V^{\pi}(s)$. Beide sehen also dieselbe Eingabe und
unterscheiden sich allein in dem, was sie ausgeben.

Damit ist zugleich gesagt, welcher Art die Agenten dieser Arbeit sind:
Angreifer wie Verteidiger sind Deep-\ac{RL}-Agenten nach genau diesem Muster,
jeder mit einem Actor- und einem Critic-Netz.

\subsection{Aktualisierung der Policy}
\label{sec:policyupdate}

Bislang wurde gesagt, dass die Policy im Training verändert wird, aber nicht,
wie. Dieser Abschnitt schließt die Lücke. Er beschreibt zunächst das Prinzip,
dann die Schätzung des Vorteils und schließlich den Ablauf, in dem beides
zusammenwirkt.

\paragraph{Das Prinzip: ein Schritt entlang des Gradienten.} Die Policy
ist nach \Cref{sec:funktionsapproximation} durch die Gewichte
$\theta$ des Actors bestimmt. Sie zu verändern heißt also, diese Zahlen zu
verändern. Gesucht ist dafür eine Zielgröße $L(\theta)$, die genau dann groß
ist, wenn die Policy gut ist. Der \emph{Gradient} $\nabla_{\theta} L(\theta)$
gibt für jedes einzelne Gewicht an, in welche Richtung und wie stark es diese
Zielgröße beeinflusst. Ein Lernschritt verschiebt die Gewichte ein Stück in
diese Richtung:
%
\begin{equation}
  \theta \leftarrow \theta + \alpha \, \nabla_{\theta} L(\theta).
  \label{eq:gradientenschritt}
\end{equation}
%
Der Faktor $\alpha$ heißt \emph{Lernrate} und legt fest, wie weit ein einzelner
Schritt reicht. Er ist klein gewählt, sodass sich die Policy über viele kleine
Schritte verändert statt über wenige große.

Damit bleibt die Frage, welche Zielgröße einzusetzen ist. Das
\emph{Policy-Gradienten-Theorem} gibt darauf die
Antwort~\cite[Abschn.~13.2--13.5]{suttonbarto2018}; in der für Actor-Critic-Ver\-fahren
gebräuchlichen Schätzform lautet der
Gradient~\cite[Gl.~(1)]{schulman2017ppo}
%
\begin{equation}
  \nabla_{\theta} L(\theta) = \hat{\mathbb{E}}_t\!\left[
    \nabla_{\theta} \log \pi_{\theta}(a_t \mid s_t)\, \hat{A}_t \right].
  \label{eq:policygradient}
\end{equation}
%
Darin ist $\hat{A}_t$ der geschätzte Vorteil der tatsächlich gespielten Aktion.
Der wahre Vorteil aus \Cref{eq:advantage} ist unbekannt, weil $Q^{\pi}$ und
$V^{\pi}$ es sind; geschätzt wird er aus den tatsächlich erhaltenen Rewards und
den Vorhersagen des Critic~\cite{schulman2018gae}.

Für jede beobachtete Kombination aus Zustand und Aktion wird ermittelt, wie die
Gewichte zu verändern wären, damit genau diese Aktion in genau diesem Zustand
wahrscheinlicher wird, und dieser Beitrag wird mit dem Vorteil gewichtet. War
der Vorteil positiv, verschiebt der Schritt die Gewichte so, dass die Aktion
künftig wahrscheinlicher wird; war er negativ, kehrt sich das Vorzeichen um und
sie wird unwahrscheinlicher. Genau darin besteht das Lernen.

Das Symbol $\hat{\mathbb{E}}_t$ bezeichnet dabei einen Erwartungswert, der nicht
exakt berechnet, sondern als Mittelwert über die gesammelten Beobachtungen
geschätzt wird. Gemittelt wird über alle Paare aus Zustand und Aktion, die der
Agent in den vorangegangenen Zeitschritten tatsächlich durchlaufen hat. Der
Unterschied zum Erwartungswert $\mathbb{E}$ aus \Cref{eq:v-aus-q} ist also, dass
dort über eine bekannte Verteilung summiert wird, während hier eine Stichprobe
an ihre Stelle tritt; das Dach kennzeichnet die Schätzung.

\paragraph{Der Ablauf.}
Beides zusammen ergibt eine Schleife, die sich über das gesamte Training
wiederholt und in \Cref{fig:rl-schleife} dargestellt ist. Der Agent handelt
zunächst eine feste Zahl von Schritten nach seiner aktuellen Policy und
speichert dabei Zustände, Aktionen, Rewards und Critic-Schätzungen; dieser
Vorrat heißt \emph{Rollout}. Anschließend werden daraus die Vorteile bestimmt.
Das setzt voraus, dass bereits feststeht, wie es weiterging, und ist deshalb
erst nach dem Sammeln möglich. Erst dann folgen die Gewichtsanpassungen: Der Vorrat wird in
kleinere Teile, sogenannte \emph{Minibatches}, zerlegt und mehrfach
durchlaufen, wobei jeder Teil einen Schritt nach \Cref{eq:gradientenschritt}
auslöst. Ein vollständiger Durchgang durch den Vorrat heißt \emph{Epoche}.
Zugleich wird der Critic daraufhin angepasst, dass seine Schätzungen den
tatsächlich beobachteten Returns näher kommen. Danach wird der Vorrat verworfen
und die Schleife beginnt von vorn, denn er stammt nun von einer veralteten
Policy.

\begin{figure}[H]
  \centering
  \begin{tikzpicture}[
    node distance=5mm,
    schritt/.style={draw=topoLine, line width=0.4pt, fill=topoNode,
      rounded corners=2pt, align=center, inner sep=4pt,
      text width=92mm, minimum height=10mm, font=\small},
    ->, >={Stealth[length=4pt]},
  ]
    \node[schritt] (a) {\textbf{1. Erfahrung sammeln}\\[1pt]
      {\scriptsize $n$ Schritte nach der aktuellen Policy handeln (Rollout)}};
    \node[schritt, below=of a] (b) {\textbf{2. Vorteile schätzen}\\[1pt]
      {\scriptsize $\hat{A}_t$ aus Rewards und Critic-Schätzungen}};
    \node[schritt, below=of b] (c) {\textbf{3. Gewichte anpassen}\\[1pt]
      {\scriptsize mehrere Epochen über Minibatches: Actor nach
      \Cref{eq:gradientenschritt}, Critic gegen die Returns}};
    \node[schritt, below=of c] (d) {\textbf{4. Erfahrung verwerfen}\\[1pt]
      {\scriptsize der Vorrat stammt jetzt von einer veralteten Policy}};
    % Ruecklaufkante rechts an allen Kaesten vorbei: Die Spalte (rand) liegt
    % ausserhalb der breitesten Box, damit der Pfeil keinen Kasten kreuzt.
    \coordinate (rand) at ([xshift=6mm]a.east);
    \draw (a) -- (b);
    \draw (b) -- (c);
    \draw (c) -- (d);
    \draw (d.east) -- (d.east -| rand) -- (a.east -| rand) -- (a.east);
  \end{tikzpicture}
  \caption[Trainingsschleife eines Actor-Critic-Verfahrens]{Trainingsschleife
  eines Actor-Critic-Verfahrens. Gehandelt wird in Schritt~1, gelernt erst in
  Schritt~3; zwischen zwei Lernphasen bleibt die Policy unverändert.}
  \label{fig:rl-schleife}
\end{figure}

Zwei Eigenschaften dieses Ablaufs sind später von Bedeutung. Erstens ist die
Policy zwischen zwei Lernphasen konstant, der Agent lernt also nicht nach jedem
einzelnen Schritt. Zweitens wird derselbe Vorrat mehrfach verwendet, was die
Zahl der benötigten Umgebungsschritte senkt, zugleich aber die Frage aufwirft,
wie weit sich die Policy dabei von derjenigen entfernen darf, mit der der Vorrat
gesammelt wurde. Genau daran setzt das im folgenden Abschnitt beschriebene
Verfahren an.

\subsection{Proximal Policy Optimization}
\label{sec:ppo}

\ac{A2C} und \ac{PPO} sind zwei verbreitete Verfahren, die beide der
Actor-Critic-Architektur folgen und regelmäßig miteinander verglichen
werden~\cite{delafuente2024comparative}. Sie unterscheiden sich darin, wie stark
ein einzelner Lernschritt die Policy verändern darf.

A2C sammelt wenige Schritte Erfahrung und aktualisiert die Policy daraufhin
unmittelbar~\cite[Abschn.~13.5]{suttonbarto2018}; die Schrittweite ist dabei
nicht nach oben begrenzt. Fällt ein Vorteil einmal ungewöhnlich groß aus, kann
eine einzelne Aktualisierung die Policy weit verschieben. Eine Aktion, die
zufällig einmal sehr gut ausging, wird dann sofort zur Gewohnheit, auch wenn sie
es nicht verdient hat. Schulman et al. benennen genau das als Schwäche des
unbeschränkten Vorgehens, das zu \glqq destructively large policy
updates\grqq{} führen könne~\cite[Abschn.~2.1]{schulman2017ppo}. \ac{PPO} setzt
dort an und deckelt die Verschiebung.

Gemessen wird sie an einer einzigen Größe: Um welchen Faktor ändert sich die
Wahrscheinlichkeit der Aktion, die tatsächlich gespielt wurde? Hielt die alte
Policy sie mit 20 Prozent für richtig und die neue mit 24 Prozent, so beträgt
dieser Faktor $24/20 = 1{,}2$. Formal ist das~\cite[Abschn.~3]{schulman2017ppo}
%
\begin{equation}
  \label{eq:ppo-ratio}
  r_t(\theta) = \frac{\pi_{\theta}(a_t \mid s_t)}{\pi_{\theta_{\text{alt}}}(a_t \mid s_t)}.
\end{equation}
%
Dabei steht $\theta$ für die Gewichte des Actor-Netzes, also für das, was beim
Training verändert wird, und $\theta_{\text{alt}}$ für ihren Stand vor der
Aktualisierung. Im Zähler steht die neue Wahrscheinlichkeit der Aktion $a_t$ im
Zustand $s_t$, im Nenner die alte. Ein Wert von eins bedeutet keine Änderung,
ein größerer Wert eine Erhöhung, ein kleinerer eine Verringerung.

Zu beachten ist, dass $r_t(\theta)$ keine einzelne Zahl je Lernschritt ist. Der
Index $t$ steht für einen bestimmten Zeitschritt, und das Verhältnis wird für
jedes im Rollout beobachtete Paar $(s_t, a_t)$ gesondert gebildet. Erst die
Mittelung über all diese Paare ergibt die Größe, die das Training maximiert.

Um diesen Faktor legt \ac{PPO} ein festes Band. Erlaubt ist nur der Bereich
zwischen $1-\varepsilon$ und $1+\varepsilon$, bei dem von Schulman et al.
vorgeschlagenen Wert $\varepsilon = 0{,}2$ also zwischen 0,8 und 1,2, was einer
Änderung von höchstens zwanzig Prozent
entspricht~\cite[Abschn.~3]{schulman2017ppo}. Verlässt der Faktor dieses Band,
wird er für die Berechnung auf den Randwert zurückgesetzt. Genau das leistet die
Funktion $\operatorname{clip}$ in der
Zielfunktion~\cite[Gl.~(7)]{schulman2017ppo}
%
\begin{equation}
  \label{eq:ppo-clip}
  L^{\text{CLIP}}(\theta) = \hat{\mathbb{E}}_t\!\left[
    \min\!\left(
      r_t(\theta)\,\hat{A}_t,\;
      \operatorname{clip}\!\left(r_t(\theta),\, 1-\varepsilon,\, 1+\varepsilon\right)\hat{A}_t
    \right)\right],
\end{equation}
%
die anstelle der Grundform aus \Cref{eq:policygradient} in
\Cref{eq:gradientenschritt} eingesetzt wird. Sie enthält drei Bestandteile. Der
geschätzte Vorteil $\hat{A}_t$ gibt an, ob die Aktion besser oder schlechter
ausfiel als erwartet. Der Faktor $r_t(\theta)$ davor sorgt dafür, dass eine
Erhöhung der Wahrscheinlichkeit bei positivem Vorteil den Ausdruck wachsen
lässt: Gute Aktionen wahrscheinlicher zu machen, lohnt sich. Das
$\hat{\mathbb{E}}_t$ ist der bereits eingeführte geschätzte Erwartungswert, hier
also der Mittelwert über alle Paare aus Zustand und Aktion des Rollouts.

Der Nutzen der Deckelung zeigt sich am Zahlenbeispiel. Eine Aktion sei mit
$\hat{A}_t = 5$ deutlich besser ausgefallen als erwartet, und die Aktualisierung
wolle ihre Wahrscheinlichkeit von 20 auf 40 Prozent anheben, also
$r_t(\theta) = 2$. Der linke Term im Minimum ergäbe dann $2 \cdot 5 = 10$, der
rechte dagegen nur $1{,}2 \cdot 5 = 6$, weil der Faktor dort auf 1,2 gestutzt
wird. Das Minimum wählt den kleineren Wert, hier also 6. Das ist genau der Wert,
den auch eine Anhebung auf lediglich 24 Prozent erbracht hätte. Über das Band
hinaus ist somit nichts mehr zu gewinnen, und der Anreiz zum großen Sprung
entfällt. Bei negativem Vorteil greift die Begrenzung spiegelbildlich auf der
Gegenseite und verhindert, dass die Wahrscheinlichkeit einer Aktion in einem Zug
fast auf null gedrückt wird.

Damit ist zugleich die am Ende von \Cref{sec:policyupdate} offene Frage
beantwortet: Der Vorrat darf mehrfach verwendet werden, weil die Deckelung
verhindert, dass sich die Policy dabei zu weit von derjenigen entfernt, mit der
er gesammelt wurde.

\subsection{Typische Schwierigkeiten}
\label{sec:rl-schwierigkeiten}

Fünf Eigenschaften des Ansatzes erschweren seinen Einsatz und werden im weiteren
Verlauf der Arbeit wiederholt eine Rolle spielen.

\paragraph{Hoher Bedarf an Erfahrung.}
Aus dem Credit-Assignment-Problem aus \Cref{sec:credit} folgt unmittelbar, dass
ein Agent sehr viele Schritte benötigt, bevor sich eine Bewertung festigt. Wo
überwachte Verfahren aus einem einzelnen eingestuften Beispiel unmittelbar
lernen, muss ein \ac{RL}-Agent denselben Zusammenhang über viele Episoden hinweg
wiederholt beobachten. Ein Training in der Realität scheidet dadurch für viele
Aufgaben aus.

\paragraph{Der Reward bestimmt das Verhalten vollständig.}
Ein Agent maximiert die Zielgröße, die ihm vorgegeben wird, und nicht die
Absicht dahinter. Belohnt eine Reward-Funktion versehentlich einen Vorgang, der
sich beliebig wiederholen lässt, ohne dass sich an der Lage etwas ändert, so
lernt der Agent genau diese Wiederholung. Sein Verhalten ist dann formal korrekt
und inhaltlich wertlos.

Damit ist die Reward-Hypothese aus \Cref{sec:rl-grundprinzip} nicht widerlegt,
aber ihre praktische Grenze benannt: Sie behauptet, dass sich ein Ziel als
skalare Größe ausdrücken \emph{lässt}, nicht, dass eine beliebige solche Größe
das gemeinte Ziel auch trifft. Ob beides zusammenfällt, entscheidet der Entwurf
der Reward-Funktion. Er ist deshalb kein Nebenaspekt, sondern bestimmt, was
überhaupt gemessen wird.

\paragraph{Mehrere lernende Agenten verletzen die Modellannahme.}
Der \ac{MDP} aus \Cref{sec:mdp} setzt eine konstante Übergangsfunktion voraus.
Handeln zwei Agenten im selben System und lernen beide gleichzeitig, so ist das
nicht mehr gegeben: Aus Sicht des einen verändert sich die Umgebung schon
dadurch, dass der andere seine Policy anpasst. Formal liegt dann kein \ac{MDP}
mehr vor, sondern ein \emph{Markov-Spiel}, und die Zielgröße verschiebt sich
fortlaufend. Aufgaben dieser Art sind Gegenstand des \ac{MARL}, das hier nicht
weiter verfolgt wird: Der Aufbau dieser Arbeit vermeidet die Schwierigkeit,
indem zu jedem Zeitpunkt nur einer der beiden Agenten lernt
(\Cref{sec:phasen}).

\paragraph{Unvollständige Beobachtbarkeit.}
Die Markov-Annahme verlangt, dass der beobachtete Zustand alle für die Zukunft
wesentliche Information enthält. In sicherheitsnahen Anwendungen ist das selten
erfüllt: Ein Angreifer kennt nur die Teile des Netzes, die er bereits entdeckt
hat. Ein Verteidiger sieht umgekehrt zwar den Zustand der eigenen Systeme und
damit auch, welche davon kompromittiert sind, nicht aber den Wissensstand des
Angreifers, also welche Knoten dieser bereits entdeckt und welche Zugangsdaten
er erbeutet hat. Beide beobachten den Zustand somit nur teilweise; formal liegt
dann ein \emph{teilweise beobachtbarer} Entscheidungsprozess vor.

In der Praxis wird die Abweichung überwiegend hingenommen und weiterhin ein
\ac{MDP} unterstellt. Gangupantulu et al. benennen den Grund: Die teilweise
beobachtbare Modellierung sei zwar realistischer, habe sich aber für große
Netze nicht als skalierbar erwiesen und verlange, zahlreiche
Wahrscheinlichkeitsverteilungen vorab zu
modellieren~\cite[Abschn.~II.C]{gangupantulu2021cja}. Die Verfahren werden also
angewendet, ohne dass ihre Voraussetzungen streng erfüllt sind, und das ist
eine bewusste Abwägung zwischen Wirklichkeitstreue und Handhabbarkeit.

\paragraph{Stochastik der Ergebnisse.}
Schließlich verläuft das Training selbst nicht deterministisch. Die
Initialisierung der Netze, die Ziehung der Aktionen aus der Policy und die
Umgebung selbst enthalten Zufall. Zwei Läufe desselben Aufbaus können deshalb
unterschiedlich ausfallen, weshalb ein einzelner Lauf eine Stichprobe darstellt
und kein Ergebnis.

\section{Grundbegriffe der Netzwerksicherheit}
\label{sec:netzsicherheit}

Dieser Abschnitt wechselt vom Lernverfahren zum Gegenstand. Er beschreibt
die Rollen von Angreifer und Verteidiger sowie den Ablauf eines Angriffs,
legt fest, was in dieser Arbeit unter einer Topologie verstanden wird, und
stellt die Segmentierungsmuster vor, aus denen die verglichenen Topologien
hervorgehen.

\subsection{Angreifer, Verteidiger und der Ablauf eines Angriffs}
\label{sec:angreifer-verteidiger}

Ein Netz besteht in der hier verwendeten Abstraktion aus \emph{Knoten}. Ein
Knoten steht für einen einzelnen Rechner, also einen \emph{Host}, der
\emph{Dienste} anbietet. Jeder Dienst lauscht auf einem \emph{Port} und verlangt
zur Anmeldung \emph{Zugangsdaten}. Welche Verbindungen zwischen zwei Knoten
überhaupt zustande kommen dürfen, regeln \emph{Firewall-Regeln}: Eine Freigabe
erlaubt Verkehr auf einem bestimmten Port, eine Sperre unterbindet ihn.

Der \emph{Angreifer} verfolgt das Ziel, Kontrolle über möglichst wertvolle
Systeme zu erlangen. Für diese wertvollsten Systeme hat sich der Begriff
\emph{Crown Jewels} eingebürgert. Gemeint sind diejenigen IT-Systeme, auf deren
bestimmungsgemäßen Betrieb eine Organisation angewiesen ist; wie kritisch ein
System ist, bemisst sich dabei an der Funktion, die es erfüllt, und an den
Daten, die es vorhält~\cite[Abschn.~II.A]{gangupantulu2021cja}.

Entscheidend an dieser Definition ist, dass sie nicht an der Netzstruktur
ansetzt, sondern an der Aufgabe der Organisation. Ermittelt werden Crown Jewels
in einem eigenen Verfahren, der \emph{Crown Jewels Analysis}, die die
Geschäftsaufgabe mit den vorhandenen IT-Werten abgleicht und jedem System einen
Relevanzwert zuweist; erst dadurch treten die betreffenden Systeme als die
wichtigsten und wertvollsten hervor~\cite[Abschn.~II.A]{gangupantulu2021cja}.
Derselbe Rechner kann somit in einer Organisation ein Crown Jewel sein und in
einer anderen ohne besondere Bedeutung.

Welches System diese Rolle in einem Netz einnimmt, ist damit keine Frage der
Technik, sondern eine Festlegung, die vorab getroffen werden muss.

Üblicherweise wird nicht der Weg in das Netz hinein modelliert, sondern der Weg
darin. Das \emph{Assumed-Breach-Prinzip} setzt voraus, dass ein erster Host
bereits kompromittiert ist, und fragt, was der Angreifer von dort aus erreicht.
Die Annahme ist realistisch, weil sich der erste Zugang, etwa über eine
Phishing-Nachricht, kaum zuverlässig verhindern lässt, und sie verschiebt den
Blick auf das, was danach geschieht.

Dieses Weiterarbeiten von einem übernommenen Host zum nächsten heißt
\emph{laterale Bewegung} und ist als eigene Taktik
definiert~\cite{mitreattack}. Sie setzt zweierlei voraus: einen Netzpfad zum
Ziel, den die Firewall zulassen muss, und gültige Zugangsdaten für das dort
laufende System. Beide Voraussetzungen sind voneinander unabhängig und müssen
gemeinsam erfüllt sein. Ein Angreifer, der ein Passwort
kennt, den betreffenden Rechner aber nicht erreicht, kommt nicht weiter.
Umgekehrt gilt dasselbe.

Ein Host, der einem Angreifer als Zwischenstation für solche Schritte dient,
heißt \emph{Pivot Point}, ein Knoten, den jeder Weg zu einem Ziel passieren
muss, \emph{Chokepoint}~\cite[Abschn.~II.A]{gangupantulu2021cja}. Ein Pivot
Point ist nicht zwingend selbst wertvoll; wird sein Risiko allein an seiner
eigenen Funktion bemessen, bleibt der Weg, den er eröffnet, unentdeckt.

Eine besondere Rolle nimmt dabei der \emph{Domain-Controller} ein, also der
Verzeichnisdienst, an dem sich die Systeme einer Windows-Domäne anmelden. Seine
Domänendatenbank enthält die Passwort-Hashes sämtlicher Domänenkonten, und deren
Auslesen ist als eigene Angriffstechnik dokumentiert~\cite{mitreattack_ntds}.
Ein Hash ist dabei noch kein Passwort: Er lässt sich offline durch
systematisches Ausprobieren einem Passwort zuordnen oder, je nach
Anmeldeverfahren, unmittelbar an dessen Stelle vorlegen. Wer den
Domain-Controller übernimmt, verfügt damit über die Anmeldeinformationen des
gesamten Netzes.

Der \emph{Verteidiger} verfolgt das Gegenziel und verfügt dafür über Maßnahmen,
die sich in zwei Gruppen einteilen lassen. \emph{Wiederherstellende} Maßnahmen
machen eine bereits eingetretene Kompromittierung rückgängig; das
\emph{Neuaufsetzen} eines Hosts, auch \emph{Reimaging} genannt, entzieht dem
Angreifer den Zugang zu diesem System. \emph{Präventive} Maßnahmen verhindern
einen Zugriff, bevor er stattfindet, etwa indem der eingehende Verkehr zu einem
gefährdeten System gesperrt wird.

Weder wiederherstellende noch präventive Maßnahmen sind kostenlos, und darin
liegt der eigentliche Zielkonflikt der Verteidigung. Ein neu aufgesetzter Host steht während des Vorgangs nicht zur
Verfügung, ein gesperrter Dienst ist auch für berechtigte Nutzer nicht mehr
erreichbar. Wie viel Ausfall hinnehmbar ist, wird als \emph{Verfügbarkeit}
gemessen und üblicherweise in einem \ac{SLA} festgeschrieben. Ein Verteidiger,
der das Netz vollständig abschaltet, hat zwar jeden Angriff unterbunden, aber
seine Aufgabe verfehlt.

\subsection{Netzwerktopologie als Sicherheitsarchitektur}
\label{sec:topologien}

Der Begriff Topologie wird uneinheitlich verwendet und bedarf daher einer
Festlegung. In dieser Arbeit bezeichnet er die Sicherheitsarchitektur eines
Netzes, also den Verlauf der Vertrauensgrenzen und die Frage, welche Systeme
miteinander kommunizieren dürfen. Gemeint ist ausdrücklich nicht die physische
Verkabelung: Graphentheoretische Muster wie Bus, Ring oder Stern beschreiben die
Verdrahtung von Geräten und liegen auf einer anderen Abstraktionsebene.

Der Grund für diese Festlegung liegt darin, dass die Verkabelung nichts darüber
aussagt, wer worauf zugreifen darf. Zwei Systeme am selben Switch können durch
Zugriffsregeln vollständig voneinander getrennt sein, während zwei Systeme an
entgegengesetzten Enden eines Gebäudes uneingeschränkt miteinander
kommunizieren. Für einen Angreifer, der sich lateral bewegt, ist allein
maßgeblich, welche Zugriffe zugelassen sind.

Die Zero-Trust-Architektur des \ac{NIST} verlagert den Maßstab in dieselbe
Richtung, von der Lage eines Systems auf die für es geltenden Regeln.
Ihr zweiter Grundsatz lautet, dass die Position im Netz für sich genommen kein
Vertrauen begründet; die zugehörigen Annahmen halten fest, dass das gesamte
private Unternehmensnetz nicht als implizite Vertrauenszone gelten
darf~\cite[Abschn.~2.1--2.2]{nist800207}. Über einen Zugriff entscheidet danach
nicht die Zugehörigkeit zu einem Netzabschnitt, sondern eine ausdrückliche
Regel. Damit ist die Ebene benannt, auf der diese Arbeit den Topologiebegriff
ansiedelt: die der zugelassenen Kommunikationsbeziehungen.

Das Mittel dazu ist die \emph{Segmentierung}. Sie unterbindet nicht benötigte
Kommunikationspfade und erschwert damit genau die laterale Bewegung aus
\Cref{sec:angreifer-verteidiger}: Ein Angreifer, der einen Host übernommen hat,
erreicht von dort aus weniger Systeme, als er es in einem unsegmentierten Netz
täte. Genau darin unterscheiden sich die im folgenden Abschnitt beschriebenen
Muster.

\subsection{Segmentierungsmuster}
\label{sec:segmentierungsmuster}

Die gebräuchlichen Muster lassen sich als Spektrum auffassen, das vom
unsegmentierten Netz bis zur feingranularen Trennung einzelner Systeme reicht.
Vier davon werden in dieser Arbeit verglichen; ihre konkrete Umsetzung
beschreibt \Cref{sec:topomodell}.

\paragraph{Flach.}
Ein flaches Netz kennt keine interne Segmentierung. Alle Hosts liegen in
derselben Vertrauenszone, jeder erreicht jeden. Das ist der Ausgangszustand, den
die Zero-Trust-Literatur als Risiko benennt: Ein einziger kompromittierter Host
genügt, um sich frei im gesamten Netz zu bewegen~\cite{nist800207}.

\paragraph{DMZ.}
Die \ac{DMZ} führt eine grobe Zonentrennung ein. Von außen erreichbare Dienste
stehen in einer entmilitarisierten Zone, die per Firewall vom internen Netz
getrennt ist~\cite{nist800041}. Der Perimeter ist damit befestigt, das Innere
bleibt aber vergleichsweise offen. Für einen Angreifer bedeutet das eine Hürde
beim Übertritt in das Intranet, danach jedoch vergleichsweise freie Bewegung.

\paragraph{Hub-and-Spoke.}
Das Hub-and-Spoke-Modell fällt aus der Reihe, weil es nicht mit Zonen arbeitet,
sondern sämtliche Zugriffe über einen einzigen Knoten führt. Als
Referenzarchitektur ist es dadurch gekennzeichnet, dass die Verbindungen nicht
transitiv sind: Jeder Spoke ist mit dem Hub verbunden, kein Spoke mit einem
anderen. Die Arbeitslasten der Spokes sind damit voneinander getrennt, während
die zentrale Zugriffskontrolle im Hub liegt~\cite{mshubspoke}. Der Hub ist
damit ein Chokepoint im Sinne von \Cref{sec:angreifer-verteidiger}, und zwar
ein baulich erzwungener: Er entsteht nicht dadurch, dass ein Angreifer diesen
Weg wählt, sondern dadurch, dass kein anderer existiert.

Der Begriff selbst stammt aus der Praxis der Netzarchitektur und ist in keinem
Standard normiert, das zugrunde liegende Prinzip ist es sehr wohl: \ac{NIST}
verlangt für eine Zero-Trust-Architektur, dass Ressourcen ohne Passieren eines
Kontrollpunkts nicht erreichbar sind~\cite[Abschn.~3.4.1]{nist800207}.
Hub-and-Spoke ist die bauliche Umsetzung dieses Gedankens: ein einziger Punkt,
an dem jeder Weg vorbeiführt.

Der Preis des Musters ist eine ausgeprägte Abhängigkeit, denn die Trennung der
Spokes ruht vollständig auf diesem einen Punkt. Geht die Kontrolle über ihn
verloren, entfällt sie ersatzlos. Realisiert wird der Hub deshalb üblicherweise
nicht durch ein Endsystem, sondern durch ein eigens dafür vorgesehenes und
entsprechend gehärtetes Gerät.

\paragraph{Micro-Segmentierung.}
Micro-Segmentierung treibt die Trennung auf die Ebene einzelner Systeme. Erlaubt
ist nur, was explizit nötig ist, und jede Verbindung wird verifiziert. Das
entspricht dem Zero-Trust-Prinzip~\cite{nist800207}. Über die Erreichbarkeit
hinaus betrifft der Grundsatz auch die Zugangsdaten: Zugriffe sollen je Sitzung
geprüft werden, statt dauerhaft gültige Berechtigungen auf den Systemen
vorzuhalten~\cite[Abschn.~2.1]{nist800207}. Beides zusammen verkürzt die Wege
eines Angreifers nicht, sondern verlängert sie, weil er für jeden Schritt sowohl
einen zugelassenen Pfad als auch die passenden Zugangsdaten benötigt.

\section{Stand der Forschung}
\label{sec:sota}

Dieser Abschnitt ordnet die Arbeit in die vorhandene Literatur ein. Er
stellt die verfügbaren Simulationsumgebungen für Netzwerkangriffe vor,
prüft, inwieweit sie einen lernenden Verteidiger vorsehen, und beschreibt
mit MARLon die Erweiterung, auf der der eigene Aufbau beruht.

\subsection{Simulationsumgebungen für Netzwerkangriffe}
\label{sec:umgebungen}

Ein \ac{RL}-Agent lernt aus Erfahrung und benötigt dafür nach
\Cref{sec:rl-schwierigkeiten} sehr viele Episoden. In einem produktiven
Firmennetz ist das weder praktikabel noch verantwortbar. Das Training findet
deshalb in einer eigens gebauten Umgebung statt, die die für laterale Bewegung
wesentlichen Mechanismen nachbildet, ohne ein reales Netz in allen Details zu
kopieren. Einen Überblick über \ac{RL}-Verfahren in der Cybersicherheit
insgesamt geben Nguyen und Reddi~\cite{nguyen2023drl}.

Solche Umgebungen unterscheiden sich vor allem entlang einer Achse. Eine
\emph{Simulation} bildet das Netz als abstraktes Modell im Programm nach; ein
Schritt kostet dann nur die Rechenzeit für eine Zustandsänderung. Eine
\emph{Emulation} betreibt stattdessen echte Systeme, etwa virtuelle Maschinen
mit realer Netzanbindung, sodass ein Schritt tatsächlich ausgeführt wird. Die
Emulation ist wirklichkeitsnäher, die Simulation um Größenordnungen schneller.

\paragraph{CyberBattleSim.}
\ac{CBS}~\cite{cyberbattlesim2021} ist eine Simulation von Microsoft. Sie bildet
ein Netzwerk als gerichteten Graphen ab; jeder Knoten steht für einen Host mit
Diensten, Schwachstellen und einer Firewall-Konfiguration. Der Angreifer beginnt
nach dem Assumed-Breach-Prinzip auf einem bereits kompromittierten
Einstiegsknoten. Eine Aktion ist etwa das Ausnutzen einer lokalen oder
entfernten Schwachstelle oder der Verbindungsaufbau zu einem weiteren Knoten mit
zuvor erbeuteten Zugangsdaten. Das Netz wird unmittelbar im Programmcode als
Graph definiert, sodass sich Knoten, Zugangsdaten und Firewall-Regeln gezielt
setzen lassen. Bereitgestellt wird die Umgebung über die
\emph{Gym-Schnittstelle}~\cite{cyberbattlesim2021}, die festlegt, wie eine
Umgebung ihren Zustand zurückgibt und eine Aktion entgegennimmt. Auf derselben
Schnittstelle setzt auch \ac{SB3} auf~\cite{stablebaselines3}, sodass sich beide
ohne Anpassung verbinden lassen.

\paragraph{CybORG.}
CybORG~\cite{standen2021cyborg} verfolgt einen zweigleisigen Ansatz und stellt
sowohl einen Simulations- als auch einen Emulationsmodus bereit, wobei letzterer
über eine Cloud-Infrastruktur betrieben wird. Das dokumentierte Szenario umfasst
drei Hosts in zwei Subnetzen. Die Umgebung ist zugleich der Ausgangspunkt der
CAGE-Wettbewerbe und damit vor allem auf vergleichbare, vorgegebene
Aufgabenstellungen ausgelegt.

\paragraph{CyGIL.}
CyGIL~\cite{li2021cygil} ist eine Emulation, die Aktionen auf tatsächlich
betriebenen virtuellen Maschinen ausführt. Die dokumentierten Szenarien umfassen
vier beziehungsweise neun Hosts. Der Preis der Wirklichkeitstreue ist die
Laufzeit: Nach eigener Messung dauert eine Episode etwa 30 Sekunden im kleineren
und 20 bis 30 Minuten im größeren Szenario, das Training eines einzelnen Agenten
entsprechend Stunden bis Tage. Eine spätere Ausbaustufe ergänzt mit CyGIL-S
einen Simulator, der aus Daten des Emulators erzeugt wird und dessen Betrieb
damit voraussetzt~\cite{li2023cygilsim}.

\paragraph{Angriffsgraphen als Gegenentwurf.}
Ein zweiter Zweig der Literatur verzichtet ganz auf eine ausführbare
Nachbildung des Netzes und fasst stattdessen einen vorab erzeugten
\emph{Angriffsgraphen} unmittelbar als \ac{MDP} auf, dessen
Übergangswahrscheinlichkeiten aus Bewertungskennzahlen bekannter Schwachstellen
stammen~\cite[Abschn.~II.C]{gangupantulu2021cja}. Der Ansatz skaliert auf
Graphen mit mehreren tausend Knoten, setzt aber voraus, dass sich der Verlauf
eines Angriffs vorab vollständig niederschreiben lässt, und kennt keinen
zweiten, gleichzeitig handelnden Agenten.

Welche dieser Umgebungen für die vorliegende Arbeit gewählt wurde und aus
welchen Gründen, behandelt \Cref{sec:umgebungswahl}.

\subsection{Lernende Agenten auf beiden Seiten}
\label{sec:beidseitig}

Ein Merkmal zieht sich durch die genannten Arbeiten: Trainiert wird
üblicherweise der Angreifer. Der Verteidiger ist, sofern überhaupt vorhanden,
eine feste regelbasierte Komponente.

CybORG erprobt in seiner Ursprungsfassung einen einzigen angreifenden Agenten
und führt die für einen lernenden Verteidiger nötigen Aktionen, Sensoren und
Aktuatoren als ersten Punkt der künftigen Arbeiten
auf~\cite{standen2021cyborg}. CyGIL ist ausdrücklich nur für das Training von
Angreifern erprobt~\cite{li2021cygil}; auch die spätere Ausbaustufe sieht zwar
begrifflich einen verteidigenden Agenten vor, trainiert in sämtlichen
Experimenten aber weiterhin allein den angreifenden~\cite{li2023cygilsim}. Dieselbe Einschränkung gilt für \ac{CBS}
selbst, wie die Autoren von CybORG in ihrem Vergleich feststellen: Ein fester
Verteidiger sei möglich, dessen Training jedoch nicht
vorgesehen~\cite{standen2021cyborg}.

Für eine Frage nach der \emph{Lerneffizienz eines Verteidigers} ist das die
entscheidende Lücke, denn sie lässt sich ohne einen zweiten lernenden Agenten
nicht stellen. Zugleich führt ein solcher Aufbau in die in
\Cref{sec:rl-schwierigkeiten} genannte Schwierigkeit: Lernen beide Seiten
gleichzeitig, liegt kein \ac{MDP} mehr vor.

\subsection{MARLon}
\label{sec:cbs}

MARLon~\cite{marlon} schließt genau diese Lücke für \ac{CBS}. Es erweitert
die Umgebung um einen zweiten, ebenfalls lernenden Agenten für die Verteidigung;
beide handeln abwechselnd im selben Netz. Bereitgestellt werden dafür
Implementierungen von \ac{A2C} und \ac{PPO} aus der Bibliothek
\ac{SB3}~\cite{stablebaselines3}, jeweils für Angreifer wie Verteidiger.

Der Verteidiger erhält dabei eine an den Angreifer gekoppelte Bewertung, und
sein Aktionsraum umfasst Maßnahmen wie das Neuaufsetzen eines Knotens sowie das
Sperren und Freigeben von Verkehr. Die genaue Ausgestaltung dieser Kopplung, die
sich für die vorliegende Arbeit als folgenreich erwiesen hat, behandelt
\Cref{sec:marlon-original}.

MARLon ist damit die Grundlage der vorliegenden Arbeit. Welche Anpassungen daran
nötig waren, beschreibt \Cref{ch:versuchsaufbau}.
